% === Begin Label Data ===
% ID: 175
% 难度星级: 
% 标签: 
% 备注: 
% 组卷引用次数: 0
% === End  Label Data ===

\begin{problem}{2015}{G}{上海卷（理）}{22}{数列}
已知数列 $\{a_n\}$ 与 $\{b_n\}$ 满足 $a_{n+1}-a_n=2(b_{n+1}-b_n), n \in \mathbf{N}^*$.

（1）若 $b_n=3n+5$, 且 $a_1=1$, 求数列 $\{a_n\}$ 的通项公式

（2）设 $\{a_n\}$ 的第 $n_0$ 项是最大项, 即 $a_{n_0} \ge a_n (n \in \mathbf{N}^*)$, 求证：数列 $\{b_n\}$ 的第 $n_0$ 项是最大项

（3）设 $a_1=\lambda<0, b_n=\lambda^n (n \in \mathbf{N}^*)$, 求 $\lambda$ 的取值范围, 使得 $\{a_n\}$ 有最大值 $M$ 与最小值 $m$, 且 $\dfrac{M}{m} \in (-2, 2)$
\end{problem}